\section{Evaluation Design}
\label{sec:evaluation}
This section specifies a proposed confirmatory study and distinguishes it from
available exploratory artifacts. The final sample size, model snapshots, budgets,
reference policy, and split have not been approved or frozen. Writing the protocol
does not mean these experiments have been run. The design emphasizes interpretable
contrasts rather than a single overall leaderboard.

\subsection{Research questions}
\begin{description}
\item[RQ1: Baseline limitations.] Which obligations do direct LLM predictions
recognize, supply, and satisfy, and which failures are hidden by coarse labels?
\item[RQ2: Context and localization.] How do location cues and surrounding
behavior affect candidate identification and mapping quality?
\item[RQ3: Same-model improvement.] Does source analysis plus semantic feedback
improve checked claim correctness and obligation completion over direct prediction
and self-review using the same model?
\item[RQ4: Mechanism.] Which benefits depend on source facts, explicit obligation
tracking, counterexample content, or simply additional inference budget?
\item[RQ5: Trust and transfer.] Are the evaluators discriminating and
tolerant of equivalent representations, and do any gains transfer to unseen source lineages and
behavior-preserving variants?
\item[RQ6: Cost and practical limits.] What model, analysis, checking, review,
and fallback costs accompany any gain, and when should the method abstain?
\end{description}

\subsection{Population, grouping, and leakage control}
The current corpus has already informed prompts, checks, and failure analysis.
It is therefore development-exposed. Its unrun input views are not an independent
holdout because they share a mother case with observed responses. The proposed
confirmatory population comprises newly reviewed source lineages, selected for
scope and contract coverage before inspecting the evaluated models' answers.
Corpus sizes and group counts remain blank in Table~\ref{tab:population}.

All views, renamings, parameter changes, and contract variants of one mother case
stay in the same partition. Shared upstream tasks and substantially shared kernels
form larger lineage groups. The split uses the larger dependency group where
appropriate. Selection and exclusions are recorded with reasons; instances are
not retained because they happen to separate models. Test generation and checker
changes after seeing responses are explicitly post-hoc diagnostics, never
retroactively part of the confirmatory score.

\begin{table}[tb]
\caption{Corpus and review report to complete after the population is fixed.
Blank cells mean unfilled, not zero or an observed unknown.}
\label{tab:population}
\begin{tabularx}{\linewidth}{Yrrrr}
\toprule
Partition & Mothers & Lineages & Views & Reviewed \\
\midrule
Development-exposed & \blank & \blank & \blank & \blank \\
Method validation & \blank & \blank & \blank & \blank \\
Confirmatory holdout & \blank & \blank & \blank & \blank \\
Search stratum & \blank & \blank & \blank & \blank \\
Optimization stratum & \blank & \blank & \blank & \blank \\
Scope / uncertain controls & \blank & \blank & \blank & \blank \\
\bottomrule
\end{tabularx}
\end{table}

The model runtime receives a frozen input allowlist and no private references,
prior responses, or final tests. For a tool-assisted run, accessible tool outputs
are separately specified, and arbitrary repository access is prohibited. A
fresh output directory, immutable raw responses, hashes, and event logs support
replay. Hashes establish artifact identity rather than truth or evaluator security.
Public pretraining contamination cannot be ruled out by local access restrictions;
source exposure and robustness analyses address it only partially.

\subsection{Experimental conditions and budget controls}
The principal unit of comparison is a fixed model configuration on the same
mother case. Direct prediction is compared with self-review, analysis-assisted
revision, semantic-feedback revision, and their combination
(Table~\ref{tab:arms}). The self-review baseline is motivated by work on iterative
self-feedback~\cite{selfrefine}, but its exact prompt and budget are frozen for
this task rather than assumed to reproduce a published implementation.

For the main revision experiment, independent initial drafts are saved and cloned
into the revision arms. Each branch receives the same initial answer, cannot see
other branches, and is allowed the same number of revision calls. Arm order is
balanced within each model and time block. A separate pre-generation analysis
contrast measures whether facts help the \emph{initial} proposal; it is not mixed
with the shared-initial-draft repair comparison.

\begin{table}[tb]
\caption{Proposed within-model conditions. All revision arms share the initial
draft in the paired study. A common contract-oriented prompt is used throughout.}
\label{tab:arms}
\begin{tabularx}{\linewidth}{lYYY}
\toprule
Arm & Source analysis & External semantic evidence & Purpose \\
\midrule
D: Direct & None & None & One-shot reference \\
S: Self-review & None & None & Extra-call control \\
A: Analysis & Source facts & None & Fact contribution \\
V: Verification & None & Development witness or finite-pass scope & Feedback contribution \\
AV: Combined & Source facts & Development witness or finite-pass scope & Proposed integration \\
\bottomrule
\end{tabularx}
\end{table}

Equal calls are not equal computation. We record prompt and completion tokens,
reasoning tokens when available, cache use, wall time, tool time, and monetary cost
when meaningful. The primary configuration uses matched call caps and declared
output limits. Sensitivity analyses use a predeclared total token cap and a wall-time
cap; realized usage remains visible. Padding prompts with irrelevant text would
not establish equivalent useful computation. Different providers' effort labels
are not treated as equivalent budgets.

Models are chosen to assess whether the mechanism generalizes beyond one model,
with exact identifiers, access path, date, revision when exposed, and configuration
recorded. Unknown temperature, seed, or server snapshot stays unknown. Existing
multi-provider pilot runs have different scaffolding and budgets and are reported
as exploratory material, not a matched capability ranking. The new model count,
replication count, timeout, and total call budget are deliberately unfilled:
models \field; replicates per mother \field; timeout \field; total calls \field.

\subsection{Outcome definitions}
We report a vector of outcomes rather than a weighted aggregate. Delivery status
distinguishes a valid final answer, schema failure, output-budget truncation,
infrastructure failure, and not-run slot. All planned slots remain in the
denominator inventory. A truncated answer supplies no successful plan, but is not
invented evidence of a false mathematical claim.

\paragraph{Localization and recognition.}
Exact source-region agreement and merged-line intersection over union are
descriptive localization diagnostics. The final scientific interpretation of
alternative boundaries uses the frozen reference policy. Structural eligibility,
intent, family, support, practical suitability, and adoption are recorded
separately. For unresolved reference labels, report coverage and abstention
without treating uncertainty as a negative class. If a stratum has no reviewed
negative references, do not report negative-class accuracy for it.

\paragraph{Claim correctness and obligation completion.}
For an applicable obligation $o$ in case $i$, let $p_{io}$ indicate supplied
content and let $v_{io}$ take supported, refuted, or unresolved values. With the
reference obligation set $O_i$ fixed independently of the answer, report
\begin{equation}
 C_i=\frac{\sum_{o\in O_i}p_{io}}{|O_i|},\qquad
 S_i=\frac{\sum_{o\in O_i}\mathbf{1}[v_{io}=\text{supported}]}{|O_i|},\qquad
 F_i=\frac{\sum_{o\in O_i}\mathbf{1}[v_{io}=\text{refuted}]}{|O_i|}.
 \label{eq:metrics}
\end{equation}
Also report the unresolved share and the evidence scope for each obligation.
These are proposed descriptive coverage measures, not a replacement for the
existing evaluator or a whole-plan correctness score. Multiple assertions within
one obligation retain separate evidence records; one supported assertion cannot
erase a contradiction or discharge unrelated requirements. The obligation-level
aggregation rule must be frozen before use. A conditional accuracy among checked
claims is accompanied by its checking coverage to prevent selective verification
from appearing comprehensive.

\paragraph{Revision transitions.}
On each applicable checked claim or explicitly scoped local task, retain the
before/after transition: refuted to supported, supported to refuted, unresolved
to supported, unresolved to refuted, or unchanged. A claim removed from the answer
may reduce false assertions while decreasing completion; it is not automatically
a successful repair. Repair yield is reported on genuinely refuted initial drafts
and alongside the full-population change. If no initial error occurs, conditional
repair capability is unmeasured rather than perfect.

\paragraph{Context obligations and fallback.}
Report validation order, exact witness or optimum selection, no-solution handling,
decoding, side effects, and caller-visible trajectories separately. A statement
that a classical fallback exists is not an executed integration test. Where a
reviewed implementation is available, report guard coverage, fallback frequency,
and cost. Otherwise these remain plan obligations. Whole-program contract
preservation and quantum resource feasibility are never imputed from local checks.

\subsection{Experiments and ablations}
\label{sec:experiments}
\paragraph{E1: Direct prediction and failure characterization (RQ1).}
Run the unhinted-context condition with first attempts preserved. Characterize
delivery, region nomination, structural judgments, planning coverage, semantic
refutations, and incomplete obligations. Compare coarse label agreement with
content-level evidence on the same responses. Code failures by mechanism and
report uncertainty rather than forcing all cases into a single error class.

\paragraph{E2: Context and location cues (RQ2).}
Run paired A/B/C views with new sessions and a balanced order. Compare B/C for
location assistance, and compare shared core obligations across A/B. Report any
change in candidate over-nomination on scope controls. Stratify descriptively by
dependency distance, branch structure, and ordering obligations. The data may
show no additional difficulty; that is a valid result.

\paragraph{E3: Paired revision (RQ3).}
Compare D, S, A, V, and AV using shared drafts for the revision arms. Measure
supported obligations, refuted claims, unresolved obligations, regression, and
completion under reserved checks. The main planned contrasts are AV versus S
and V versus S, with A versus S and AV versus V isolating the added analysis.
Select the confirmatory contrast before running the study; report all remaining
contrasts as secondary. Include correct initial drafts so that regressions and
ceiling effects remain visible.

\paragraph{E4: Component and information ablations (RQ4).}
Remove source facts, remove the obligation ledger while retaining task content,
replace a witness with a verdict-only message, and omit the violated-obligation
explanation. Each comparison changes one factor within its parent condition.
Add the separate pre-generation fact condition and budget-matched self-review.
No arm receives final-test answers. Assess whether richer feedback helps repair,
merely elicits longer plans, or causes unjustified confidence. All ablations are
proposals pending implementation and authorization.

\paragraph{E5: Evaluator validity and transcription (RQ5).}
Test independent task definitions against known-correct, equivalent, and faulty
constructions. Cover variable permutation, bit complement, constant energy shifts,
optimization-preserving scaling, and invalid-state extrema. Audit model-passage
to claim transcription independently, including ambiguous formulas. Measure
unsupported representations, disagreement, transcription time, and false acceptance
or rejection on controls whose expected behavior is independently established.
Never report a control mutation kill rate as model improvement.

\paragraph{E6: Reachability and behavioral robustness (RQ5).}
Vary input order, tied optima, empty/no-solution cases, invalid inputs, numerical
boundaries, and caller budgets where the public contract allows them. Distinguish
behavior-preserving transformations from contract-changing variants. For numerical
counterexamples, provide a path from valid inputs to the tested intermediate
state. Compare checking a local expression alone with checking it under the
source-state conditions. These tests do not add unsupported quantum families.

\paragraph{E7: Transfer and model sensitivity (RQ5).}
Freeze the method on development lineages and evaluate unseen lineages without
updating prompts, facts, or checkers from their responses. Report within-model
gains separately for search, optimization, and unresolved controls. Repeat the
same assistance contrasts across models, preserving provider-specific settings.
Leave-one-lineage-out development analysis can guide uncertainty estimates, but
it is not relabeled as an untouched final holdout.

\paragraph{E8: Efficiency and adoption limits (RQ6).}
Report model calls, tokens, wall time, analysis time, exact-checking time, peak
memory when measured, human review effort, and fallback cost. Plot quality against
the recorded call/token/time budgets only after results exist. Vary bounded
enumeration size to expose checker scaling. Preserve timeout and unsupported
outcomes. If circuit resource data are collected, report them separately with
representation and backend assumptions; no hardware advantage is inferred.

\subsection{Statistical analysis and stopping rules}
The independent sampling unit is the mother problem or, when shared provenance
requires it, the larger lineage group. Repeated drafts, context views, obligations,
and individual checker instances are nested observations. For a per-mother metric
$y$ and paired arms $a,b$, first average within the fixed replicate set to obtain
$d_i=\overline{y}_{i,a}-\overline{y}_{i,b}$. Report the distribution of $d_i$ and
the lineage-weighted contrast, giving each independent lineage equal weight.
Do not multiply a formula's test instances by response count to enlarge sample size.

The proposed interval procedure resamples whole independent lineage groups while
retaining all paired arms and nested observations. The confidence level, bootstrap
replications, minimum useful effect, and required interval precision are fixed
before execution. If groups are too few for reliable interval inference, report
the paired observations descriptively and state that limitation. A power or
precision analysis may use development variability, but must not promise that a
convenient case count is sufficient. Confirmatory multiplicity correction applies
to the predeclared family of contrasts; exploratory strata remain labeled as such.

Missing delivery contributes to the fixed-slot outcome inventory and cannot be
removed to improve reported success. Analyses restricted to checkable responses
state their reduced coverage. No mathematical correctness label is imputed for
missing output. Collection stops at the fixed budget or a documented infrastructure
stop rule. There are no selective retries, extra repetitions after disappointing
results, or model substitutions within the registered experiment.

\subsection{Reproducibility and research use of AI}
The replication record should include protocol and source hashes, raw provider
responses, timing and usage, tool events, parsing decisions, claim transcriptions,
review provenance, environment locks, and derived-table scripts. Development
feedback and reserved evaluation inputs are stored separately. Anonymous reviewer
access must be tested before submission; this draft does not claim a deposited
artifact.

AI tools have assisted development of case proposals, code, checking infrastructure,
review drafts, and this methodological draft. Models are also the experimental
subjects. These roles are distinct: an AI-proposed reference is not an independent
human annotation, and a generated checker is not trusted solely because it runs.
The final methods record must identify the tools, versions where available,
artifacts affected, human oversight, and validation for each research role. No
independent expert agreement is inferred from an unattributed review response.
